\chapter{Spec-Drift Audit Agent}
\label{ch:spec-drift-agent}

\section{Purpose and scope}

The Spec-Drift Audit Agent is a specialized governance agent that monitors
synchronization between three artifact types:

\begin{enumerate}
  \item \textbf{LaTeX specifications} in \filepath{docs/latex/specs/}
  \item \textbf{JSON/YAML schemas} in \filepath{docs/schema/}
  \item \textbf{Implementation code} in \filepath{src/}
\end{enumerate}

When any of these artifacts change, the agent ensures related artifacts are
also updated to maintain consistency. This prevents \emph{specification drift}
where documentation diverges from implementation.

\section{Agent configuration files}

\subsection{Primary configuration}

\begin{table}[htbp]
\centering
\footnotesize
\caption{Spec-drift agent configuration files.}
\label{tab:spec-drift-files}
\begin{tabularx}{\textwidth}{@{}l X@{}}
\toprule
\textbf{File} & \textbf{Purpose} \\
\midrule
\filepath{.clinerules.spec-drift-auditor} & Agent behavior rules and drift type definitions \\
\filepath{docs/schema/spec-code-mapping.yaml} & Traceability registry mapping artifacts \\
\filepath{docs/schema/drift-exceptions.yaml} & Approved exceptions with expiration \\
\filepath{scripts/spec-drift/audit.py} & Drift detection engine implementation \\
\filepath{scripts/spec-drift/pre-commit-hook.sh} & Git pre-commit hook \\
\filepath{.github/workflows/spec-drift-audit.yml} & CI/CD workflow \\
\bottomrule
\end{tabularx}
\end{table}

\subsection{Domain-specific rule packs}

Each specification directory contains a \filepath{.clinerules} file that
defines domain-specific synchronization rules:

\begin{table}[htbp]
\centering
\footnotesize
\caption{Domain-specific .clinerules locations.}
\label{tab:domain-clinerules}
\begin{tabularx}{\textwidth}{@{}l X@{}}
\toprule
\textbf{Location} & \textbf{Domain} \\
\midrule
\filepath{docs/schema/.clinerules} & Schema governance and data contracts \\
\filepath{docs/latex/specs/.clinerules} & Specification directory standards \\
\filepath{docs/latex/specs/simula/.clinerules} & Training data framework \\
\filepath{docs/latex/specs/regulations/.clinerules} & AI regulatory compliance \\
\filepath{docs/latex/specs/arabic/.clinerules} & Arabic AP Invoice processing \\
\filepath{docs/latex/specs/tb/.clinerules} & Trial Balance review \\
\filepath{docs/latex/specs/tb-hitl/.clinerules} & TB business requirements \\
\filepath{docs/latex/specs/arabic-hitl/.clinerules} & Arabic business requirements \\
\filepath{docs/latex/specs/clinerules-agents/.clinerules} & This specification \\
\bottomrule
\end{tabularx}
\end{table}

\section{Drift types detected}
\label{sec:drift-types}

The audit agent detects eight categories of specification drift:

\begin{table}[htbp]
\centering
\footnotesize
\caption{Drift type classification.}
\label{tab:drift-types}
\begin{tabularx}{\textwidth}{@{}l l X@{}}
\toprule
\textbf{Code} & \textbf{Severity} & \textbf{Description} \\
\midrule
DRIFT-001 & HIGH/MEDIUM & Schema-Spec Drift: Schema field changed without spec update \\
DRIFT-002 & CRITICAL/HIGH & Code-Schema Drift: Code references non-existent schema fields \\
DRIFT-003 & CRITICAL/HIGH & Code-Spec Drift: Code behavior contradicts spec requirement \\
DRIFT-004 & HIGH/MEDIUM & Version Drift: Version numbers out of sync \\
DRIFT-005 & CRITICAL/HIGH & Cross-Domain Drift: Shared enum changed without updating consumers \\
DRIFT-006 & HIGH/MEDIUM & Threshold Drift: Code uses different threshold than spec \\
DRIFT-007 & CRITICAL/HIGH & State Machine Drift: Workflow states differ between code and spec \\
DRIFT-008 & CRITICAL/HIGH & API Contract Drift: API differs from spec documentation \\
\bottomrule
\end{tabularx}
\end{table}

\section{Enforcement mechanisms}
\label{sec:enforcement}

\textbf{CRITICAL: To guarantee spec-drift checks always run, multiple
enforcement layers are required.} No single mechanism is sufficient.

\subsection{Layer 1: CI/CD enforcement (Mandatory)}

The GitHub Actions workflow at \filepath{.github/workflows/spec-drift-audit.yml}
provides server-side enforcement that \textbf{cannot be bypassed} by developers:

\begin{lstlisting}[language=yamlx,caption=CI/CD enforcement configuration]
on:
  pull_request:
    branches: [main, develop]
  push:
    branches: [main]

jobs:
  spec-drift-audit:
    runs-on: ubuntu-latest
    steps:
      - uses: actions/checkout@v4
      - name: Run spec-drift audit
        run: make audit-spec-drift-ci
      # Fails PR if CRITICAL or HIGH drift detected
\end{lstlisting}

This layer is \textbf{mandatory} because:
\begin{itemize}
  \item It runs on the server, not developer machines
  \item Branch protection rules can require it to pass
  \item Developers cannot skip or disable it locally
  \item All PRs are audited regardless of local setup
\end{itemize}

\subsection{Layer 2: Local pre-commit hook (Recommended)}

The pre-commit hook provides early feedback during development:

\begin{lstlisting}[language=bash,caption=Pre-commit hook installation]
# One-time setup per developer
make audit-install-hook
\end{lstlisting}

This layer is \textbf{recommended} but not mandatory because:
\begin{itemize}
  \item Developers must opt-in by running the install command
  \item Hooks can be bypassed with \texttt{git commit --no-verify}
  \item It improves developer experience but doesn't guarantee compliance
\end{itemize}

\subsection{Layer 3: Manual audit commands (On-demand)}

Developers can run audits manually at any time:

\begin{lstlisting}[language=bash,caption=Manual audit commands]
# Full audit
make audit-spec-drift

# Domain-specific audit
make audit-spec-drift-domain DOMAIN=simula

# Quick check on local changes
make audit-spec-drift-quick
\end{lstlisting}

\section{Developer onboarding}
\label{sec:onboarding}

New developers MUST complete the following setup:

\begin{enumerate}
  \item \textbf{Clone repository}: Standard git clone
  \item \textbf{Install Python dependencies}: \texttt{pip install pyyaml}
  \item \textbf{Install pre-commit hook}: \texttt{make audit-install-hook}
  \item \textbf{Verify installation}: \texttt{make audit-spec-drift}
\end{enumerate}

The onboarding checklist is documented in \filepath{scripts/spec-drift/README.md}.

\textbf{Note}: Even if a developer skips step 3, the CI/CD enforcement
(Layer 1) will catch any drift when they create a pull request.

\section{Traceability registry}
\label{sec:traceability}

The \filepath{docs/schema/spec-code-mapping.yaml} file defines relationships
between artifacts:

\begin{lstlisting}[language=yamlx,caption=Example traceability mapping]
domains:
  simula:
    spec_root: docs/latex/specs/simula
    schema_root: docs/schema/simula
    code_root: src/training
    
    artifacts:
      - id: simula-taxonomy
        type: chapter
        spec_path: docs/latex/specs/simula/chapters/05-taxonomy-engine.tex
        related_schemas:
          - docs/schema/simula/taxonomy.schema.json
        related_code:
          - src/training/pipeline/simula_taxonomy_builder.py
        sync_rules:
          - "Taxonomy node structure changes require spec+schema+code update"
\end{lstlisting}

When any file in a relationship changes, the audit agent flags all related
files as potentially requiring updates.

\section{Drift exceptions}
\label{sec:exceptions}

Temporary exceptions can be granted via \filepath{docs/schema/drift-exceptions.yaml}:

\begin{lstlisting}[language=yamlx,caption=Drift exception format]
exceptions:
  - id: "EXC-001"
    created: "2026-04-21"
    expires: "2026-06-21"      # REQUIRED: Max 90 days
    owner: "team-lead@example.com"
    drift_type: "DRIFT-003"
    file_pattern: "src/training/experiments/.*"
    rationale: |
      Experimental code is not yet production-ready.
    resolution_plan: |
      Formalize or remove before GA release.
    tracking_issue: "https://github.com/org/repo/issues/123"
\end{lstlisting}

\textbf{Exception governance rules:}
\begin{itemize}
  \item All exceptions MUST have an expiration date (max 90 days)
  \item Exceptions MUST have an owner and tracking issue
  \item Expired exceptions are flagged in audit reports
  \item Permanent exceptions require Architecture Review Board approval
\end{itemize}

\section{Minimum acceptance criteria}

The spec-drift agent MUST satisfy:

\begin{enumerate}
  \item \textbf{CI/CD workflow exists} and is configured as a required status check
  \item \textbf{Traceability registry covers} all domain specifications
  \item \textbf{CRITICAL/HIGH drift} blocks PR merges
  \item \textbf{Audit reports} are archived for 90 days
  \item \textbf{Documentation} includes onboarding and troubleshooting guides
  \item \textbf{All .clinerules files} reference the traceability registry
\end{enumerate}

\section{Guaranteeing enforcement}
\label{sec:guarantee}

To \textbf{guarantee} spec-drift checks always run:

\begin{enumerate}
  \item \textbf{Enable branch protection} on \texttt{main} and \texttt{develop}:
    \begin{itemize}
      \item Require status checks to pass before merging
      \item Add \texttt{spec-drift-audit} as a required status check
      \item Disable force pushes and branch deletion
    \end{itemize}
  
  \item \textbf{Configure repository settings}:
    \begin{itemize}
      \item Disable direct commits to protected branches
      \item Require pull request reviews
      \item Require signed commits (optional but recommended)
    \end{itemize}
  
  \item \textbf{Monitor compliance} via audit logs:
    \begin{itemize}
      \item GitHub Actions logs all audit runs
      \item Failed audits are visible in PR status
      \item Administrators can review bypass attempts
    \end{itemize}
\end{enumerate}

\begin{tcolorbox}[colback=yellow!5,colframe=yellow!70,title=\faIcon{exclamation-triangle} Enforcement Configuration Required]
The GitHub repository MUST have branch protection rules configured with
\texttt{spec-drift-audit} as a required status check. Without this
configuration, developers can merge PRs that bypass drift detection.

\textbf{Action required}: Repository administrators must configure branch
protection rules in GitHub settings $\rightarrow$ Branches $\rightarrow$ Add rule.
\end{tcolorbox}